\section{Normal length near a divisor}
\label{sec:divisorial-normal-length}

To prove completeness near a divisor, one needs a lower bound for the metric in the normal direction. A lower bound for the determinant alone does not give it, because the tangential eigenvalues may absorb the volume. Dividing the determinant by the determinant of the tangential block gives the normal Schur complement. A lower bound for it by the inverse square of a defining function of the divisor, weakened by at most the square of the logarithm of that function, forces completeness. For solutions of exponential equations, this bound follows from bounds on the density, the potential and the tangential determinant, and an example on a torus shows that the logarithmic exponent $2$ is sharp.

In this section $X$ is again a connected compact complex manifold of dimension $n\geq1$ with a Hermitian form $\omega$. Let $D\subset X$ be a nonempty smooth effective divisor, and put $M=X\setminus D$. Let $s_D$ be the canonical section of $\mathcal O_X(D)$, equipped with a smooth Hermitian norm, and set
\[
 R=|s_D|,\qquad \tau=-\log R^2.
\]
Rescale the norm so that $\tau\geq1$ on $M$. Write $J$ for the complex structure. For a smooth positive real $(1,1)$-form $T$ on $M$, we use the Riemannian metric $g_T(v,v)=\pi T(v,Jv)$ on real tangent vectors $v$. Thus, if $T=(\mathrm{i}/2\pi)\sum H_{j\bar k}\,\mathrm{d}z_j\wedge\mathrm{d}\bar z_k$, then $g_T(v,v)=\sum H_{j\bar k}\,\mathrm{d}z_j(v)\overline{\mathrm{d}z_k(v)}$, and Hermitian norms below use this normalization of coefficients.

\Needspace{7\baselineskip}
Near $D$, the form $\partial\tau$ does not vanish (see \eqref{eq:normal-gradient-comparison} below), and we define the tangential hyperplane field, the tangential block and the determinant
\[
 E_\tau=\ker\partial\tau,\qquad B_\tau=T|_{E_\tau},\qquad
 \det_\omega T=\frac{T^n}{\omega^n}.
\]
Similarly, $\det_\omega B_\tau$ denotes the determinant of $B_\tau$ relative to the restriction of $\omega$ to the complex hyperplane $E_\tau$; when $n=1$, it is $1$.

\begin{proposition}\label{prop:divisorial-normal-length}\index{Metric completeness!normal length}
Suppose that there are constants $c>0$, $\tau_0\geq1$ and $p\leq2$ such that, on $\{\tau>\tau_0\}$,
\begin{equation}\label{eq:normal-determinant-ratio}
 \frac{\det_\omega T}{\det_\omega B_\tau}
 \geq cR^{-2}\tau^{-p}.
\end{equation}
Then $(M,g_T)$ is complete.
\end{proposition}

\begin{proof}
In a coordinate neighborhood meeting $D$, write $s_D=z_1 e$, where $e$ is a local holomorphic frame with $|e|^2=\mathrm{e}^{-\psi}$ and $\psi$ is smooth. Then
\[
 \tau=-\log|z_1|^2+\psi,\qquad
 \partial\tau=-\frac{\mathrm{d}z_1}{z_1}+\partial\psi.
\]
Covering the compact divisor by finitely many such neighborhoods and comparing with the smooth background, we obtain, after increasing $\tau_0$,
\begin{equation}\label{eq:normal-gradient-comparison}
 c_1R^{-2}\leq |\partial\tau|_\omega^2\leq C_1R^{-2}
 \quad\hbox{on }\{\tau>\tau_0\}
\end{equation}
for positive constants $c_1,C_1$. In particular, $E_\tau$ is a hyperplane there.

At a point of this set, choose an $\omega$-unitary frame whose last $n-1$ vectors span $E_\tau$. Let $H$ be the Hermitian matrix that represents $g_T$ as the quadratic form $w^*Hw$ in the complex components $w$, and write
\[
 H=\begin{pmatrix}h&a^*\\a&B\end{pmatrix},\qquad
 S=h-a^*B^{-1}a=\frac{\det H}{\det B},
\]
where $a$ is a column vector and $a^*$ its conjugate transpose. Completing the square gives, for a vector with components $(\xi,\eta)$,
\[
 \begin{pmatrix}\bar\xi&\eta^*\end{pmatrix}
 H\begin{pmatrix}\xi\\\eta\end{pmatrix}
 =S|\xi|^2+(\eta+B^{-1}a\xi)^*B(\eta+B^{-1}a\xi)
 \geq S|\xi|^2.
\]
For $n=1$ the same identity holds without the second term. For the complex components of a real vector $v$ we have
$|\xi|=|\partial\tau(v)|/|\partial\tau|_\omega$, so \eqref{eq:normal-determinant-ratio} and \eqref{eq:normal-gradient-comparison} imply
\begin{equation}\label{eq:normal-length-bound}
 g_T(v,v)\geq c_2\tau^{-p}|\partial\tau(v)|^2
 \geq \frac{c_2}{4}\tau^{-p}|\mathrm{d}\tau(v)|^2,
\end{equation}
where $c_2>0$ is uniform. The last inequality uses $\mathrm{d}\tau(v)=2\operatorname{Re}\partial\tau(v)$.

Define a continuous function $\Phi_p$ on $[1,\infty)$ by
\[
 \Phi_p(t)=
 \begin{cases}
  0,&t\leq\tau_0,\\
  \displaystyle\int_{\tau_0}^t x^{-p/2}\,\mathrm{d}x,&t>\tau_0.
 \end{cases}
\]
Integrating \eqref{eq:normal-length-bound} along piecewise smooth curves shows that $\Phi_p\circ\tau$ is Lipschitz for the length distance of $g_T$; this also holds for curves crossing $\{\tau=\tau_0\}$, since $\Phi_p$ is locally Lipschitz and constant below that level. Since $p\leq2$, $\Phi_p(t)$ tends to infinity with $t$, so $\Phi_p\circ\tau$ is a proper function on $M$.

A Cauchy sequence for the length distance has bounded values of this Lipschitz function and hence lies in a compact subset of $M$. It therefore has a subsequence converging in the manifold topology. On a coordinate ball about the limit, the smooth positive form $T$ is comparable to the Euclidean metric, so this subsequence also converges for the length distance. The Cauchy property then gives convergence of the whole sequence, which proves completeness.
\end{proof}

\Needspace{16\baselineskip}
\begin{corollary}\label{cor:divisorial-exponential-completeness}\index{Exponential equation!normal completeness}
In the preceding setting, let $\theta$ be a smooth real $(1,1)$-form on $X$, and let $u\in C^\infty(M)$ satisfy $T=\theta+\ddc u>0$ on $M$. Suppose that, for some $\lambda>0$ and some positive $f\in C^\infty(M)$,
\[
 T^n=\mathrm{e}^{\lambda u}f\omega^n\quad\hbox{on }M.
\]
Assume that, near $D$, there are constants $c,C>0$ and $q,a,b\in\mathbb R$ such that
\begin{equation}\label{eq:divisorial-three-bounds}
 f\geq cR^{-2}\tau^{-q},\qquad
 u\geq-\frac{a}{\lambda}\log\tau-C,\qquad
 \det_\omega B_\tau\leq C\tau^b.
\end{equation}
If $q+a+b\leq2$, then $(M,g_T)$ is complete.
\end{corollary}

\begin{proof}
The equation and the first two inequalities give
$\det_\omega T\geq c\mathrm{e}^{-\lambda C}R^{-2}\tau^{-(q+a)}$.
Dividing by the third bound gives \eqref{eq:normal-determinant-ratio} with $p=q+a+b$, and \cref{prop:divisorial-normal-length} applies.
\end{proof}

The relevant tangential determinant is that of the intrinsic hyperplane field $E_\tau$, not that of an arbitrary block of coordinate directions.

\subsection{The critical exponent on a compact torus}

The exponent $2$ in \cref{cor:divisorial-exponential-completeness} is sharp. On a product torus, two explicit logarithmic profiles give a complete metric with $q+a+b=2$ and incomplete metrics with $q+a+b=2+\beta$ for every $\beta>0$. Both profiles have full mass, and the second one is bounded.

Let $E=\mathbb C/(\mathbb Z+\mathrm{i}\mathbb Z)$, let $\omega_E=(\mathrm{i}/2\pi)\,\mathrm{d}z\wedge\mathrm{d}\bar z$, and put $X=E^n$, $D=\{0\}\times E^{n-1}$ and $M=X\setminus D$. Write $\pi_j\colon X\to E$ for the coordinate projections and $\omega_0=\sum_{j=1}^n\pi_j^*\omega_E$. Choose radii $0<r_1<r_2<\mathrm{e}^{-1}$ and a smooth radial cutoff $\chi$ on the coordinate disk $\{|z|<r_2\}\subset E$, equal to $1$ on $\{|z|\leq r_1\}$ and to $0$ near its outer boundary, with $0\leq\chi\leq1$. For fixed numbers $A,\eta,\beta>0$, set $L=-\log|z|^2$ and define on $E$
\begin{equation}\label{eq:torus-length-profiles}
 v_{\mathrm c}=-A\chi\log L,\qquad
 v_\beta=\eta\chi L^{-\beta}.
\end{equation}
Both functions are extended by zero outside the disk, and at its center we set $v_{\mathrm c}(0)=-\infty$ and $v_\beta(0)=0$. The bounded profile also occurs in \cref{ex:logarithmic-continuous-density}. We write $\mathrm{MA}_\theta(u)$ for the extended non-pluripolar Monge--Amp\`ere measure $\mathrm{MA}^{\mathrm{ext}}_\theta(u)$ of \cref{thm:arbitrary-extended-ma}.

\Needspace{14\baselineskip}
\begin{lemma}\label{lem:torus-logarithmic-full-mass}
For $C>0$ sufficiently large, put $\theta=C\omega_0$. For either profile $v$ in \eqref{eq:torus-length-profiles}, the function $u=\pi_1^*v$ belongs to $\PSH(X,\theta)$, the form $T=\theta+\ddc u$ is smooth and positive on $X\setminus D$, and
\[
 \int_X\mathrm{MA}_\theta(u)=\int_X\theta^n.
\]
The same $C$ may be used for the two fixed profiles. The bounded profile has minimal ordinary singularities; the profile $v_{\mathrm c}$ does not.
\end{lemma}

\begin{proof}
For a twice differentiable function $G$ of $L=-\log|z|^2$, differentiation away from zero gives
\[
 \frac{\partial^2G(L)}{\partial z\partial\bar z}
 =\frac{G''(L)}{|z|^2}.
\]
On the inner disk this yields
\begin{equation}\label{eq:torus-profile-hessians}
 (v_{\mathrm c})_{z\bar z}=\frac{A}{|z|^2L^2},\qquad
 (v_\beta)_{z\bar z}=\frac{\eta\beta(\beta+1)}{|z|^2L^{\beta+2}}.
\end{equation}
Both densities are integrable at zero. Both profiles are locally integrable, and their distributional Hessians have no additional atom at zero. To see this, integrate twice by parts outside $\{|z|=\varepsilon\}$: the boundary terms against a smooth test function are bounded by constant multiples of
$\varepsilon|G(-\log\varepsilon^2)|$ and $|G'(-\log\varepsilon^2)|$, and both tend to zero for the two profiles. Thus \eqref{eq:torus-profile-hessians} holds in the sense of distributions across the center. On the cutoff annulus all derivatives are bounded. Hence, for a sufficiently large common $C$,
\[
 \alpha_v=C\omega_E+\ddc v
\]
is a positive measure, smooth and strictly positive away from zero, with no atom at zero. Pairing $\ddc v$ with the constant test function gives
$\int_E\alpha_v=C\int_E\omega_E$.

Let $Y=E^{n-1}$, let $\pi_Y\colon X\to Y$ be the projection onto the remaining factors, and let $\theta_Y$ be $C$ times the sum of the $n-1$ flat forms of $Y$. On $X\setminus D$ the smooth wedge product is
\[
 T^n=n\,\pi_1^*\alpha_v\wedge\pi_Y^*\theta_Y^{n-1};
\]
for $n=1$ this means just $\alpha_v$, with $Y$ a point. The right side defines a finite product measure on $X$ that gives $D$ zero mass. By \cref{thm:arbitrary-extended-ma}, the non-pluripolar measure agrees with this product off $D$ and gives the pluripolar set $D$ zero mass, so the two global measures are equal. Integrating the product proves the full-mass identity. Since $\theta>0$, its envelope is $V_\theta=0$ (\cref{sec:obstacle-envelopes}). The bounded profile therefore has minimal ordinary singularities, whereas $v_{\mathrm c}$ is not bounded below by $0$ minus a constant.
\end{proof}

We apply the preceding results with $\omega=\theta$. Choose a smooth norm on $\mathcal O_E(0)$, the line bundle of the point $0\in E$, for which its canonical section has squared norm $|z|^2$ near zero; such a norm is obtained from any smooth norm by multiplying by a smooth positive function prescribed on a smaller disk. Pull it back to $X$ and rescale it by a constant. For some fixed $K>0$ we may then arrange
\[
 R^2=\mathrm{e}^{-K}|z_1|^2,\qquad \tau=L+K
\]
near $D$, and $\tau\geq1$ on $M$. For both profiles, $E_\tau$ is exactly the tangent space of the remaining torus factors near $D$, so $\det_\theta B_\tau=1$.

\Needspace{8\baselineskip}
\begin{example}\label{ex:critical-complete-divisor}\index{Metric completeness!critical logarithmic exponent}
Fix $\lambda>0$ and $0<A<1/\lambda$, and take $u_{\mathrm c}=\pi_1^*v_{\mathrm c}$ with the background of \cref{lem:torus-logarithmic-full-mass}. This gives a complete metric on $M$ solving
\[
 \mathrm{MA}_\theta(u_{\mathrm c})
 =\mathrm{e}^{\lambda u_{\mathrm c}}f_{\mathrm c}\theta^n,
 \qquad f_{\mathrm c}\in C^\infty(M)\cap L^1(X,\theta^n),\quad f_{\mathrm c}>0,
\]
where the right side is extended by zero on $D$. The three exponents in \eqref{eq:divisorial-three-bounds} can be taken to be
\[
 q=2-\lambda A>1,\qquad a=\lambda A,\qquad b=0,
\]
so their sum equals the critical exponent $2$.

Indeed, define $f_{\mathrm c}=\mathrm{e}^{-\lambda u_{\mathrm c}}\det_\theta(\theta+\ddc u_{\mathrm c})$ on $M$. On the inner disk the product structure gives
\[
 f_{\mathrm c}=L^{\lambda A}
       +\frac{A}{C|z_1|^2L^{2-\lambda A}}.
\]
The substitution $L=-\log r^2$ in the radial integral shows that the second term is integrable precisely when $\lambda A<1$. The first term is integrable as well, and the density is smooth elsewhere. The identity of measures across $D$ follows from the preceding lemma. Since $\tau=L+K$, all three bounds hold with the stated exponents, and completeness follows from \cref{cor:divisorial-exponential-completeness}. Along a radial normal segment the divergence of the length at the critical exponent is also visible directly:
\[
 \int_0^{r_1}\frac{\sqrt A}{r(-\log r^2)}\,\mathrm{d}r=\infty.
\]
For arbitrary paths, however, completeness requires the corollary, not just this family of radial paths.
\end{example}

\Needspace{9\baselineskip}
\begin{example}\label{ex:supercritical-incomplete-divisor}
Fix $\lambda,\eta,\beta>0$ and take $u_\beta=\pi_1^*v_\beta$ with the background of \cref{lem:torus-logarithmic-full-mass}. This is a continuous bounded solution of full mass with minimal ordinary singularities, and its metric on $M$ is incomplete, although
\[
 \mathrm{MA}_\theta(u_\beta)=\mathrm{e}^{\lambda u_\beta}f_\beta\theta^n,
 \qquad f_\beta\in C^\infty(M)\cap L^1(X,\theta^n),\quad f_\beta>0.
\]
To see this, put $K_\beta=\eta\beta(\beta+1)$ and define the density by the equation on $M$. Near $D$ it is
\[
 f_\beta=\mathrm{e}^{-\lambda\eta L^{-\beta}}
       \left(1+\frac{K_\beta}{C|z_1|^2L^{\beta+2}}\right).
\]
It is integrable, and the global identity follows as before. Fix the other coordinates and let $z_1=r\downarrow0$ along a radial segment starting at $r_0$, with $0<r_0<r_1$. The length of this segment is at most
\[
 \sqrt C\,r_0+\sqrt{K_\beta}
   \int_0^{r_0}\frac{\mathrm{d}r}{r(-\log r^2)^{1+\beta/2}}
 =\sqrt C\,r_0+\frac{\sqrt{K_\beta}}{\beta}
       (-\log r_0^2)^{-\beta/2}<\infty.
\]
Points along this segment approaching $D$ form a Cauchy sequence with no limit in $M$.

Here the smallest admissible exponents in \eqref{eq:divisorial-three-bounds} are $q=\beta+2$, $a=0$ and $b=0$: the density is asymptotic to a positive constant times $R^{-2}\tau^{-(\beta+2)}$, the potential is bounded, and the tangential determinant is $1$. Their sum is $2+\beta$. Since $\beta>0$ is arbitrary, the threshold $2$ in the criterion cannot be replaced by any larger universal number, even in a K\"ahler class with a bounded full-mass solution.
\end{example}

In these examples one may write $f=F/R^2$ with $F=R^2f$ smooth and positive on $M$, but $F$ does not extend smoothly and positively across $D$: the logarithmic factors make it tend to zero there. Thus the examples concern integrable divisorial poles weakened by logarithmic factors. They do not decide completeness for a nonintegrable pole with smooth positive coefficient across the divisor, nor the question of compatibility with the flow discussed in \cref{sec:history-completeness}.
